\documentclass[11pt]{article}

\usepackage{amsmath,amssymb,amsthm,mathtools}
\usepackage{hyperref}
\usepackage{enumitem}
\usepackage{geometry}
\geometry{margin=1in}

\title{Prime-Colored Braid Multiplicity:\\
A Conditional, Testable Program for Knot Invariants}

\author{Citizen Gardens \\ \\ The Foundation of Multiplicity}

\date{April 4, 2026}

\newtheorem{theorem}{Theorem}
\newtheorem{conjecture}{Conjecture}
\newtheorem{definition}{Definition}
\newtheorem{remark}{Remark}

\begin{document}

\maketitle

\begin{abstract}
We construct a prime-colored braid representation whose local data are driven by prime arithmetic, and use it to define a conditional, numerically testable knot functional.
The construction is organized as a program: we specify the algebraic layer (prime-colored local unitaries and strand-dependent $R$-matrices), the canonical braid and symmetrization layer, the prime-harmonic parameters, and a protection functional that couples a prime-weighted trace to a topological complexity term.
We prove that the prime-harmonic parameters converge under mild regularity assumptions, and isolate three remaining conjectural layers: a projective Yang--Baxter property with vanishing residuals, a restricted Markov invariance, and numerical convergence of the protection functional.
Together these define a sharp target for both computation and theory: either the program yields a genuinely new knot invariant, or explicit numerical counterexamples reveal where it fails.
\end{abstract}

\section{Executive summary}

This paper proposes a \emph{prime-colored braid multiplicity} framework with the following structure.

\begin{itemize}[leftmargin=1.5em]
  \item \textbf{Local data.}
  To each prime $p$ we attach a unitary $O_p\in SU(2)$ depending on prime data (e.g.\ $\ln p$ and $p^{-1/2}$).
  Conjugation of a fixed $U_q(\mathfrak{sl}_2)$ $R$-matrix $R_{\mathrm{std}}$ at $q=e^{i\pi/3}$ by $O_p\otimes O_q$ yields strand-dependent $R$-matrices $R_{p,q}$.%
  \item \textbf{Prime-colored braid representation.}
  Given a braid $\beta$ on $n$ strands with prime labels $(p_1,\dots,p_n)$, we obtain a representation
  \[
    \rho(\beta;\mathbf p)\in\mathrm{End}((\mathbb C^2)^{\otimes n})
  \]
  by assigning $R_{p_i,p_{i+1}}^{\pm1}$ to crossings.
  A projective Yang--Baxter conjecture asserts that this is a projective representation of the braid group.
  \item \textbf{Canonical braid and symmetrization.}
  For a knot or link $K$, a deterministic algorithm produces a canonical braid $\beta_K$.
  We define a symmetrized trace
  \[
    Z_{\mathrm{sym}}(K;X)
    = \frac{1}{n!}\sum_{\sigma\in S_n}
      \mathrm{Tr}\bigl(\rho(\beta_K;\sigma\cdot \mathbf p)\bigr),
  \]
  where $\sigma$ permutes the base prime assignment and $X$ is a prime cutoff.
  Under natural covariance assumptions this average equals the trace of $\rho(\beta_K;\mathbf p)$ projected onto the $S_n$-invariant subspace.
  \item \textbf{Prime-harmonic parameters.}
  Using the prime-weighted measure
  \[
    \mu_p = \frac{1}{p\ln p}, \qquad
    \langle f\rangle_X = \frac{\sum_{p\le X} f(p)\,\mu_p}{\sum_{p\le X}\mu_p},
  \]
  and bounded functions $f_0,f_z$ of $\ln p$, we define
  \[
    c_0(X) = \langle f_0(p)\rangle_X,
    \quad
    z(X)   = \langle f_z(p)\rangle_X.
  \]
  Prime number theorem arguments show that $c_0(X)$ and $z(X)$ converge as $X\to\infty$.
  \item \textbf{Protection functional.}
  Choosing a knot complexity $\mathcal C(K)$ (e.g.\ crossing number or hyperbolic volume), we define
  \[
    P(K;X) = \big|Z_{\mathrm{sym}}(K;X)\big|^{z(X)}\exp\big(-c_0(X)\,\mathcal C(K)\big).
  \]
  For each finite $X$ this is computable and presentation-independent under the symmetrization.
  \item \textbf{Conditional main statement.}
  Assuming (i) projective Yang--Baxter with vanishing residuals, (ii) a suitable restricted Markov invariance, and (iii) convergence of $P(K;X)$ as $X\to\infty$, the limit
  \[
    P(K) = \lim_{X\to\infty}P(K;X)
  \]
  defines a knot invariant depending only on the isotopy class of $K$.
\end{itemize}

The remainder of the paper formalizes these ingredients, proves the convergence of the prime-harmonic parameters, and outlines a numerical protocol to validate or falsify the conjectural layers.

\section{Prime-colored braid category and local data}

\subsection{Prime-colored objects and morphisms}

Fix an infinite ordered list of primes $(p_1,p_2,p_3,\dots)$.
We define a \emph{prime-colored braid category} $\mathcal B_{\mathrm{prime}}$ as follows.

\begin{itemize}[leftmargin=1.5em]
  \item Objects: finite ordered sequences of primes $(p_{i_1},\dots,p_{i_n})$.
  \item Morphisms: braids on $n$ strands whose $j$-th strand is colored by $p_{i_j}$.
  \item Monoidal structure: tensor product given by concatenation of sequences and horizontal juxtaposition of braids.
\end{itemize}

Thus each strand carries a fixed prime label throughout its history.

\subsection{Prime-indexed local unitaries}

For each prime $p$ we choose a unitary $O_p\in SU(2)$.
A concrete choice is of the form
\[
  O_p = \exp\bigl(i\,\ln p\,(\hat n_p\cdot \sigma)\bigr),
\]
where $\sigma=(\sigma_x,\sigma_y,\sigma_z)$ are the Pauli matrices, and $\hat n_p$ is a unit vector depending on arithmetic data of $p$ (for example via the fractional part of $\alpha\ln p$ for some fixed irrational $\alpha$).\footnote{The exact form of $\hat n_p$ is modeling data, not derived; once fixed, it is used uniformly throughout.}

This choice guarantees unitarity while encoding nontrivial dependence on the prime distribution.

\subsection{Strand-dependent $R$-matrix}

Let $R_{\mathrm{std}}$ be the standard $U_q(\mathfrak{sl}_2)$ $R$-matrix in the fundamental representation at $q=e^{i\pi/3}$.
We define, for primes $p,q$,
\[
  R_{p,q} = (O_p\otimes O_q)\,R_{\mathrm{std}}\,(O_p^\dagger\otimes O_q^\dagger).
\]

Given a prime-colored braid $\beta$ on $n$ strands with color sequence $(p_1,\dots,p_n)$, we obtain a representation
\[
  \rho(\beta;\mathbf p)\in\mathrm{End}\bigl((\mathbb C^2)^{\otimes n}\bigr)
\]
by assigning $R_{p_i,p_{i+1}}$ or $R_{p_i,p_{i+1}}^{-1}$ to each elementary crossing of strands $i$ and $i+1$ according to the braid word.

\section{Projective Yang--Baxter structure}

Define
\[
  L_{p,q,r} = R_{p,q}R_{p,r}R_{q,r},\qquad
  R_{p,q,r} = R_{q,r}R_{p,r}R_{p,q}.
\]

\begin{conjecture}[Prime-labeled YBE, projective form]
\label{conj:YBE}
For all primes $p,q,r$, there exists a phase $\lambda_{p,q,r}\in U(1)$ such that
\[
  L_{p,q,r} = \lambda_{p,q,r}\,R_{p,q,r}.
\]
Moreover, the optimal residual
\[
  \epsilon_{\mathrm{YBE}}(P_{\max})
   = \max_{p,q,r\le P_{\max}}\min_{\lambda\in U(1)}
      \|L_{p,q,r}-\lambda R_{p,q,r}\|
\]
tends to $0$ as $P_{\max}\to\infty$.
\end{conjecture}

Under Conjecture~\ref{conj:YBE}, $\rho(\,\cdot\,;\mathbf p)$ is a projective representation of the braid group.
Standard 2-cocycle machinery then lifts it to a genuine representation of a central extension of the braid group, but we will work at the projective level.

\section{Canonical braid and symmetrized trace}

\subsection{Canonical braid representative}

Given an oriented diagram $D$ of a knot or link $K$, we fix a deterministic procedure:

\begin{enumerate}[leftmargin=1.5em]
  \item \emph{Planar normalization.}
  Place $D$ in a standard rectangle with orientation chosen by a deterministic rule (for example, minimizing a chosen functional of the diagram).
  \item \emph{Deterministic braid representative.}
  Apply a fixed algorithm (Alexander's theorem implemented by a prescribed sequence of moves) to produce a braid word on $n$ strands whose closure recovers $D$.
  Denote this braid by $\beta_K$.
  \item \emph{Base prime assignment.}
  Assign primes $(p_1,\dots,p_n)$ to strands $1,\dots,n$ in order.
\end{enumerate}

This yields a canonical pair $(\beta_K,\mathbf p)$ associated to $K$.

\subsection{Symmetrized prime assignment and projection}

Define the base assignment $\mathbf p=(p_1,\dots,p_n)$ and let $S_n$ act by permutation:
\[
  \sigma\cdot\mathbf p = (p_{\sigma(1)},\dots,p_{\sigma(n)}).
\]
We then define the symmetrized trace
\[
  Z_{\mathrm{sym}}(K;X)
  = \frac1{n!}\sum_{\sigma\in S_n}
    \mathrm{Tr}\bigl(\rho(\beta_K;\sigma\cdot\mathbf p)\bigr),
\]
where $X$ is a prime cutoff used implicitly in the definition of the local data and parameters.

Under natural covariance assumptions (the permutation of primes corresponds to a permutation of tensor factors), this average may be rewritten as
\[
  Z_{\mathrm{sym}}(K;X)
  = \mathrm{Tr}\bigl(\Pi_{\mathrm{triv}}\rho(\beta_K;\mathbf p)\bigr),
\]
where $\Pi_{\mathrm{triv}}$ is the projector onto the $S_n$-invariant subspace of $(\mathbb C^2)^{\otimes n}$.
This interpretation yields a more efficient implementation in practice.

\section{Prime-harmonic parameters}

\subsection{Prime-weighted measure}

Following the v3.0 framework, we define a prime-counting measure
\[
  \mu_p = \frac{1}{p\ln p},\qquad
  \langle f\rangle_X
   = \frac{\sum_{p\le X} f(p)\mu_p}{\sum_{p\le X}\mu_p}.
\]

Here $X$ is a finite cutoff; $\langle f\rangle_X$ is a truncated average over primes.

\subsection{Bounded prime-harmonic parameters}

Let $f_0,f_z:\mathbb R\to\mathbb R$ be bounded, measurable functions of $\ln p$ (and possibly of bounded functions of $p^{-1/2}$).
Set
\[
  c_0(X) = \langle f_0(p)\rangle_X,\qquad
  z(X)   = \langle f_z(p)\rangle_X.
\]

\begin{theorem}[Prime-harmonic convergence]
\label{thm:prime-limit}
Under the assumptions above on $f_0$ and $f_z$, the limits
\[
  c_0 := \lim_{X\to\infty} c_0(X),\qquad
  z   := \lim_{X\to\infty} z(X)
\]
exist and are finite.
\end{theorem}

\begin{proof}[Sketch]
This is a standard application of the prime number theorem and summation by parts.
Since $\mu_p$ approximates a probability measure on primes with density $1/(p\ln p)$, and $f_0,f_z$ are bounded and sufficiently regular in $\ln p$, the normalized sums converge to integrals of $f_0,f_z$ against an exponential kernel $e^{-u}\,du$ in the variable $u=\ln p$.
We omit the detailed analytic number theory, which follows known prime-averaging techniques.
\end{proof}

Any specific closed form for $c_0$ or $z$ (e.g.\ in terms of values of $\zeta(s)$) is an empirical hypothesis to be tested numerically; Theorem~\ref{thm:prime-limit} only asserts existence and finiteness.

\section{Protection functional}

\subsection{Complexity term}

Let $\mathcal C(K)$ be a fixed complexity measure of the knot or link $K$.
For simplicity, we may take $\mathcal C(K)$ to be the crossing number of a chosen canonical diagram, but more refined choices such as (normalized) hyperbolic volume are also possible.

\subsection{Definition and conditional invariance}

\begin{definition}[Protection functional]
For a finite prime cutoff $X$, define the protection functional
\[
  P(K;X) = \big|Z_{\mathrm{sym}}(K;X)\big|^{z(X)}\exp\big(-c_0(X)\,\mathcal C(K)\big),
\]
where $c_0(X)$ and $z(X)$ are the prime-harmonic parameters above.
\end{definition}

For each fixed $X$, $P(K;X)$ is a well-defined, computable quantity given a choice of canonical braid and complexity measure.

\begin{conjecture}[Conditional knot invariance]
\label{conj:main}
Assume:
\begin{enumerate}[label=(\roman*),leftmargin=1.5em]
  \item The projective Yang--Baxter Conjecture~\ref{conj:YBE} holds with residual $\epsilon_{\mathrm{YBE}}(P_{\max})\to 0$.
  \item A restricted Markov invariance theorem: the symmetrized functional $Z_{\mathrm{sym}}(K;X)$ is invariant under braid moves implementing Markov conjugations and stabilizations, including appropriate averaging over the prime label of the new strand.
  \item The limit $P(K)=\lim_{X\to\infty}P(K;X)$ exists for all $K$.
\end{enumerate}
Then $P(K)$ depends only on the isotopy class of $K$ and defines a knot (or link) invariant.
\end{conjecture}

Even without passing to the limit, each $P(K;X)$ is a presentation-independent statistic within the symmetrized, prime-colored braid framework.

\section{Numerical protocol and pilot code}

\subsection{Testing the projective YBE}

For a practical cutoff $P_{\max}$, we propose:

\begin{enumerate}[leftmargin=1.5em]
  \item Generate all primes $p,q,r\le P_{\max}$.
  \item For each triple, construct $O_p,O_q,O_r$ and the associated $R_{p,q}$, $R_{p,r}$, $R_{q,r}$.
  \item Compute $L_{p,q,r}$ and $R_{p,q,r}$, find the phase $\lambda$ minimizing $\|L_{p,q,r}-\lambda R_{p,q,r}\|$, and record the residual.
  \item Track $\epsilon_{\mathrm{YBE}}(P_{\max})$ as $P_{\max}$ increases.
\end{enumerate}

\subsection{Computing $Z_{\mathrm{sym}}(K;X)$}

For a small knot $K$ (e.g.\ unknot, trefoil, figure-eight):

\begin{enumerate}[leftmargin=1.5em]
  \item Construct a canonical braid $\beta_K$ on $n$ strands and the base prime assignment $\mathbf p=(p_1,\dots,p_n)$.
  \item For all $\sigma\in S_n$ (or a large random subset if $n$ is large), evaluate $\mathrm{Tr}(\rho(\beta_K;\sigma\cdot\mathbf p))$.
  \item Average over $\sigma$ to obtain $Z_{\mathrm{sym}}(K;X)$.
\end{enumerate}

\subsection{Illustrative Python-style pseudocode}

Below is a schematic snippet indicating how one might implement the core computations in Python/NumPy; it is not optimized, but illustrates the ingredients.

\begin{verbatim}
import numpy as np
from itertools import permutations

def pauli_matrices():
    sx = np.array([[0, 1], [1, 0]], dtype=complex)
    sy = np.array([[0, -1j], [1j, 0]], dtype=complex)
    sz = np.array([[1, 0], [0, -1]], dtype=complex)
    return sx, sy, sz

def O_p(p, alpha=1.0):
    sx, sy, sz = pauli_matrices()
    # Example direction from fractional part of alpha * ln p
    t = (alpha * np.log(p)) % (2*np.pi)
    nx, ny, nz = np.cos(t), np.sin(t), 0.0
    n_dot_sigma = nx*sx + ny*sy + nz*sz
    return expm(1j * np.log(p) * n_dot_sigma)

def R_std(q=np.exp(1j*np.pi/3)):
    # Hard-code the 4x4 fundamental R-matrix for U_q(sl2)
    # (omitted here for brevity)
    pass

def R_pq(p, q):
    Op = O_p(p)
    Oq = O_p(q)
    R = R_std()
    return np.kron(Op, Oq) @ R @ np.kron(Op.conj().T, Oq.conj().T)

def rho_beta(beta_word, primes):
    # beta_word: list of (i, sign) meaning sigma_i^{sign}
    # primes: list [p1,...,pn]
    dim = 2**len(primes)
    M = np.eye(dim, dtype=complex)
    for i, sgn in beta_word:
        Rloc = R_pq(primes[i], primes[i+1])
        # Embed Rloc into full tensor product (omitted)
        R_full = embed_R(Rloc, i, len(primes))
        if sgn == -1:
            R_full = np.linalg.inv(R_full)
        M = R_full @ M
    return M

def Z_sym(beta_word, primes):
    n = len(primes)
    acc = 0.0 + 0.0j
    for sigma in permutations(range(n)):
        perm_primes = [primes[j] for j in sigma]
        M = rho_beta(beta_word, perm_primes)
        acc += np.trace(M)
    return acc / np.math.factorial(n)
\end{verbatim}

A production implementation would:
(i) hard-code $R_{\mathrm{std}}$,
(ii) optimize tensor embeddings,
(iii) use the projection interpretation instead of explicit permutation averages, and
(iv) incorporate the prime-harmonic parameters and protection factor.

\section{Outlook}

This v1 framework deliberately separates what is proved (existence of prime-harmonic limits, algebraic consistency of the construction) from what is conjectural (projective YBE, Markov invariance, convergence of $P(K;X)$) and what is empirical (numerical behaviour of the residuals and invariants).
The immediate next step is a pilot numerical study for modest prime cutoffs and low-crossing knots, which will either support the conjectural layers or provide concrete counterexamples and thus sharpen the theory further.


\nocite{*}
\bibliographystyle{plainnat}
\bibliography{references}

\end{document}